\subsection{2.14 Die mathematische Strategie der PST – Vom ontologischen Fundament zur beobachtbaren Physik}

Nach Abschluss der ersten numerischen Untersuchungen lässt sich erstmals eine
klare Forschungsstrategie formulieren.

Diese Strategie unterscheidet sich grundlegend vom Vorgehen der etablierten
Physik.

Während viele heutige Theorien versuchen, bekannte physikalische Gesetze durch
weitere Annahmen zu erweitern, verfolgt die PST den entgegengesetzten Weg:

\[
\boxed{
\text{Möglichst wenige ontologische Annahmen}
\Longrightarrow
\text{möglichst viele physikalische Konsequenzen.}
}
\]

Die Aufgabe besteht daher nicht darin, bekannte Physik nachzubauen,
sondern sie als notwendige Konsequenz einer minimalen Ontologie zu rekonstruieren.

%%%%%%%%%%%%%%%%%%%%%%%%%%%%%%%%%%%%%%%%%%%%%%%%%%%%%%%%%%%%%%%

\subsubsection*{2.14.1 Die Rekonstruktionskette}

Die bisherige Entwicklung legt folgende logische Abfolge nahe:

\[
\boxed{
\text{Existenz}
\longrightarrow
\text{OPP}
\longrightarrow
\mathcal P
\longrightarrow
\text{Raumzeit}
\longrightarrow
\text{Physik}.
}
\]

Jeder Pfeil beschreibt einen Rekonstruktionsschritt.

Jeder Schritt muss mathematisch begründet werden.

Die Theorie darf an keiner Stelle zusätzliche physikalische Strukturen
einführen, wenn diese nicht aus den vorhergehenden Schritten folgen.

%%%%%%%%%%%%%%%%%%%%%%%%%%%%%%%%%%%%%%%%%%%%%%%%%%%%%%%%%%%%%%%

\subsubsection*{2.14.2 Prinzip der minimalen Ontologie}

Der OP verfolgt konsequent das Prinzip:

\[
\boxed{
\text{Ontologisch so wenig wie möglich,
mathematisch so viel wie nötig.}
}
\]

Dies bedeutet:

Der OPP soll selbst möglichst strukturfrei bleiben.

Insbesondere werden zunächst nicht vorausgesetzt:

\begin{itemize}

\item Raum,

\item Zeit,

\item Metriken,

\item Koordinaten,

\item Dimension,

\item Teilchen,

\item Felder,

\item Wechselwirkungen.

\end{itemize}

All diese Größen sollen später emergieren.

%%%%%%%%%%%%%%%%%%%%%%%%%%%%%%%%%%%%%%%%%%%%%%%%%%%%%%%%%%%%%%%

\subsubsection*{2.14.3 Warum Emergenz unvermeidlich wird}

Gerade diese ontologische Sparsamkeit erzwingt Emergenz.

Denn wenn der OPP selbst keine Raumzeit besitzt,

muss erklärt werden,

wie dennoch eine Raumzeit beobachtet werden kann.

Dasselbe gilt für

\begin{itemize}

\item Symmetrien,

\item Teilchen,

\item Massen,

\item Eichfelder,

\item Kausalität,

\item Dynamik.

\end{itemize}

Die Theorie ist daher gezwungen,

alle diese Eigenschaften aus Projektionsprozessen herzuleiten.

%%%%%%%%%%%%%%%%%%%%%%%%%%%%%%%%%%%%%%%%%%%%%%%%%%%%%%%%%%%%%%%

\subsubsection*{2.14.4 Die Rolle der Mathematik}

Hier verändert sich die Rolle der Mathematik grundlegend.

In vielen physikalischen Theorien dient Mathematik dazu,

bereits bekannte physikalische Strukturen zu beschreiben.

In der PST übernimmt Mathematik dagegen eine konstruktive Funktion.

Sie erzeugt die beobachtbare Physik überhaupt erst.

Formal:

\[
\boxed{
\text{Ontologie}
+
\text{Mathematik}
=
\text{Physik}.
}
\]

Die Mathematik ist somit nicht bloß Sprache,

sondern Teil des Erklärungsmechanismus.

%%%%%%%%%%%%%%%%%%%%%%%%%%%%%%%%%%%%%%%%%%%%%%%%%%%%%%%%%%%%%%%

\subsubsection*{2.14.5 Die Projektionsklasse}

Bisher wurden verschiedene mathematische Projektionen untersucht.

Unter anderem:

\begin{itemize}

\item PCA,

\item Kernel PCA,

\item Diffusion Maps,

\item Random Projection.

\end{itemize}

Diese Verfahren besitzen jedoch lediglich heuristischen Charakter.

Sie zeigen,

dass geeignete Projektionen tatsächlich vierdimensionale Strukturen erzeugen können.

Sie beweisen jedoch noch nicht,

welcher Operator physikalisch korrekt ist.

%%%%%%%%%%%%%%%%%%%%%%%%%%%%%%%%%%%%%%%%%%%%%%%%%%%%%%%%%%%%%%%

\subsubsection*{2.14.6 Die eigentliche Aufgabe}

Die eigentliche mathematische Herausforderung lautet daher:

\begin{center}

\Large

Nicht

\[
\boxed{\text{Welche Projektion funktioniert?}}
\]

sondern

\[
\boxed{\text{Warum besitzt die Natur genau diese Projektion?}}
\]

\end{center}

Erst wenn diese Frage beantwortet ist,

kann die PST als fundamentale Theorie gelten.

%%%%%%%%%%%%%%%%%%%%%%%%%%%%%%%%%%%%%%%%%%%%%%%%%%%%%%%%%%%%%%%

\subsubsection*{2.14.7 Kriterien für den gesuchten Operator}

Der gesuchte Projektionsoperator

\[
\mathcal P^\ast
\]

muss zahlreiche Bedingungen gleichzeitig erfüllen.

Insbesondere sollte er

\begin{enumerate}

\item eine stabile effektive Dimension

\[
d_{\mathrm{eff}}
\approx
4
\]

erzeugen,

\item Lorentz-Invarianz ermöglichen,

\item kontinuierliche Raumzeit hervorbringen,

\item Quantisierung erklären,

\item Eichsymmetrien zulassen,

\item robuste Ergebnisse gegenüber kleinen Störungen liefern.

\end{enumerate}

Diese Bedingungen bilden den zukünftigen mathematischen Prüfstein der PST.

%%%%%%%%%%%%%%%%%%%%%%%%%%%%%%%%%%%%%%%%%%%%%%%%%%%%%%%%%%%%%%%

\subsubsection*{2.14.8 Die Bedeutung negativer Resultate}

Ein bemerkenswerter Aspekt der bisherigen Arbeiten ist,

dass auch scheinbar negative Ergebnisse wertvoll sind.

Beispielsweise zeigte sich:

\begin{itemize}

\item Der rohe OPP besitzt keine Vierdimensionalität.

\item Beliebige Projektionen liefern keine stabile Physik.

\item Einige Verfahren erzeugen offensichtlich unphysikalische Ergebnisse.

\end{itemize}

Jedes dieser Resultate verkleinert den Raum möglicher Theorien.

Auch Ausschlüsse besitzen daher wissenschaftlichen Wert.

%%%%%%%%%%%%%%%%%%%%%%%%%%%%%%%%%%%%%%%%%%%%%%%%%%%%%%%%%%%%%%%

\subsubsection*{2.14.9 Rekonstruktion statt Anpassung}

Ein wesentliches Ziel der PST besteht darin,

freie Parameter möglichst zu vermeiden.

Die Theorie soll nicht an Beobachtungen angepasst werden.

Vielmehr sollen Beobachtungen aus ihr folgen.

Idealerweise ergibt sich beispielsweise:

\[
3+1
\]

nicht,

weil dies experimentell bekannt ist,

sondern weil jede andere Dimensionszahl mathematisch instabil wäre.

Dasselbe gilt später für

\begin{itemize}

\item Lichtgeschwindigkeit,

\item Eichgruppen,

\item Massenhierarchien,

\item Kopplungskonstanten.

\end{itemize}

%%%%%%%%%%%%%%%%%%%%%%%%%%%%%%%%%%%%%%%%%%%%%%%%%%%%%%%%%%%%%%%

\subsubsection*{2.14.10 Verbindung zum OP}

Diese Strategie ist unmittelbar mit dem Ontophänomenalismus verbunden.

Der OP beginnt mit der minimalen Aussage

\[
\boxed{
\text{Existenz existiert.}
}
\]

Von dort aus soll Schritt für Schritt rekonstruiert werden:

\[
\text{Existenz}
\rightarrow
\text{Bestimmtheit}
\rightarrow
\text{Differenz}
\rightarrow
\text{Relation}
\rightarrow
\text{OPP}
\rightarrow
\mathcal P
\rightarrow
\text{Raumzeit}
\rightarrow
\text{Physik}.
\]

Damit wird deutlich,

dass die PST keine eigenständige Theorie neben dem OP ist,

sondern dessen mathematische Ausarbeitung.

%%%%%%%%%%%%%%%%%%%%%%%%%%%%%%%%%%%%%%%%%%%%%%%%%%%%%%%%%%%%%%%

\subsubsection*{2.14.11 Das langfristige Ziel}

Langfristig soll die gesamte bekannte Physik

als notwendige Konsequenz einer einzigen ontologischen Grundannahme erscheinen.

Das Fernziel lautet:

\[
\boxed{
\text{Existenz}
\Longrightarrow
\text{Universum}.
}
\]

Zwischen beiden Begriffen dürfen ausschließlich mathematisch notwendige
Rekonstruktionsschritte stehen.

%%%%%%%%%%%%%%%%%%%%%%%%%%%%%%%%%%%%%%%%%%%%%%%%%%%%%%%%%%%%%%%

\subsubsection*{2.14.12 Schlussbemerkung}

Die bisherigen numerischen Untersuchungen markieren erst den Beginn dieser
Rekonstruktion.

Sie zeigen jedoch,

dass die grundlegende Forschungsrichtung tragfähig ist.

Nicht der OPP selbst scheint die entscheidende mathematische Herausforderung
darzustellen,

sondern die Struktur jener Projektion,

welche aus einem zeitlosen, hochdimensionalen ontologischen Fundament die
vierdimensionale, lorentzinvariante und quantisierte Welt der Physik entstehen
lässt.

Damit verschiebt sich der Schwerpunkt der weiteren Forschung endgültig von der
Analyse des OPP zur Herleitung des universellen Projektionsoperators
\(\mathcal P^\ast\), der als zentrales Bindeglied zwischen Ontologie und Physik
fungiert.

